The action category in the general criterion has an explicit description
directly from the coefficient diagram, without a word presentation.
\begin{theorem}[coefficient colimits]\label{thm:action-colimits}
Let $T=(Q,F)$ be any unital inverse semiring and $L$ any left
$Q$-semimodule. For $l,k\in L$, put
\[
 K(l,k)=\{q\in Q:ql\leq k\}.
\]
This is a nonempty directed poset, and there are canonical identifications
\[
 \mathcal A_L(l,k)=\varinjlim_{q\in K(l,k)}F(q).
\]
The composite of representatives $x\in F(q)$ at $l\to h$ and
$y\in F(p)$ at $h\to k$ is represented by
$\mu_{p,q}(y,x)\in F(pq)$; object identities are represented by the
global scalar unit. Let $D_l$ be the directed poset of $Q$-linear
functionals admissible at $l$, ordered pointwise. Then
\[
 \mathcal J_L(l,k)=\im\left(
  \varinjlim_{\phi\in D_l}F(\phi(k))
  \longrightarrow\varinjlim_{q\in K(l,k)}F(q)\right).
\]
Every ideal element thus has a single admissible functional and
coefficient representative. The displayed map need not be injective.

If $Q$ is finite, with no finiteness assumption on $L$ or the groups,
put $r(l,k)=\bigvee K(l,k)$. The greatest admissible functional
$\delta_l=\bigvee_{\phi\in D_l}\phi$ also exists. In that case
\[
 \mathcal A_L(l,k)=F(r(l,k)),\qquad
 \mathcal J_L(l,k)=
 \im\bigl(F(\delta_l(k))\longrightarrow F(r(l,k))\bigr).
\]
Composition first uses $\mu$ at $r(h,k)r(l,h)$ and then the coefficient
transition to $r(l,k)$. For finite $L$, these are the exact blocks and
multiplication of the ordinary category ring, even when its groups are
infinite.
\end{theorem}
The unrestricted ordinary category and its canonical ideal are therefore
known explicitly from $(Q,F)$ and the support action. Classifying their
finite projectives and the allowed projective supports remains the
residual problem; the theorem does not assert a finite algorithm for
arbitrary infinite data.
